 \documentclass[10pt]{article}
\usepackage{amsmath, amssymb,graphicx,tikz,wasysym,youngtab,comment}

\usepackage{amsfonts}
\usepackage{amsthm}
\usepackage{textcomp}
\usepackage{mdwlist}
\usepackage{enumerate}
\usepackage{multicol}
\usepackage{amsmath, array, graphicx}
\usepackage{tikz}
\usepackage{todonotes}
\usepackage{pgfplots}
\usepackage{fontawesome}
\usepackage{enumitem, multicol}
\usepackage{fancyhdr}
\usepackage{wrapfig}
\usepackage[makeroom]{cancel}
\usepackage{booktabs}
\usepackage{comment}
\usepackage{alltt}
\usepackage{pdfpages}
%\usepackage[dvips]{graphics}
\textheight9in
\textwidth17cm
\oddsidemargin0cm
\evensidemargin0cm
\setlength{\topmargin}{-0.8in}

\newcommand{\be}{\begin{enumerate}}
\newcommand{\bi}{\begin{itemize}}

\newcommand{\ei}{\end{itemize}}
\newcommand{\ee}{\end{enumerate}}
\newcommand{\ds}{\displaystyle}
\newcommand{\dx}{\, dx}
\newcommand{\ZZ}{\mathbb{Z}}
\newcommand{\QQ}{\mathbb{Q}}
\newcommand{\RR}{\mathbb{R}}
\newcommand{\CC}{\mathbb{C}}
\newcommand{\xx}{\mathbf{x}}
\newcommand{\pp}{\mathbf{p}}
%\newcommand{\null}{\mathrm{null}}
\newcommand{\row}{\mathrm{row}}
\newcommand{\col}{\mathrm{col}}
\newcommand{\nul}{\mathrm{null}}
\newcommand{\rank}{\mathrm{rank}}
\newcommand{\nullity}{\mathrm{nullity}}
\newtheorem{proposition}{Proposition}
\newtheorem{lemma}{Lemma}
\newtheorem{theorem}{Theorem}
\newtheorem{definition}{Definition}
\newtheorem{example}{Example}
\pagestyle{empty}

\newcolumntype{P}[1]{>{\centering\arraybackslash}p{#1}}
\newcolumntype{M}[1]{>{\centering\arraybackslash}m{#1}}
\newcolumntype{N}{@{}m{0pt}@{}}
\newcolumntype{L}[1]{>{\raggedright\arraybackslash}m{#1}}
\newcolumntype{R}[1]{>{\raggedleft\arraybackslash}m{#1}}
\newcommand{\babble}[1]{\marginpar{\flushleft\scriptsize\textsf{#1}}}

\oddsidemargin=0in
\textwidth=5.75in
\topmargin=-0.5in
\textheight=9in
\renewcommand{\marginparwidth}{81pt}

\begin{document}
\begin{center}
Math 364 \quad Spring 2026 \quad Problem Set \#6\\
Khanh Tra (Fay) Nguyen Tran
\end{center}

\begin{enumerate}

\item
\textbf{In class, we saw that if $\mathcal{C}_n$ is the set of lattice paths that go from $(0,0)$ to $(n,n)$ and do not go above the line $y=x$, then $|\mathcal{C}_n|= \binom{2n}{n}-\binom{2n}{n-1} = \frac{1}{n+1}\binom{2n}{n}$.}

\textbf{Now, let $\mathcal{C}_{m,n}$ be the set of lattice paths that go from $(0,0)$ to $(m,n)$ and do not go above the line $y=x$.  Find, with proof, $|\mathcal{C}_{m,n}|$.}.

Let $A_{m,n}$ be the set of lattice paths that go from $(0,0)$ to $(m,n)$ by only moving UP or RIGHT for each step. Obviously, we have $|A_{m,n}| = \binom{m+n}{n}$ since we need to take a total of $m+n$ steps in each path and in each path there are $m$ RIGHT steps and $n$ UP steps.

Let $B_{m,n}$ be the set of lattice paths that go from $(0,0)$ to $(m,n)$ by only moving UP or RIGHT for each step and touch $y = x + 1$ at least once. $B_{m,n}$ is the set that contains the bad paths from $(0,0)$ to $(m,n)$. 

We have: 

$$|\mathcal{C}_{m,n}| = |A_{m,n}| - |B_{m,n}|$$

We consider two cases:

\begin{itemize}
    \item \textbf{Case 1:} $m \geq n $

For every path $P \in B_{m,n}$, $P$ must touch the line $y=x+1$ at least once. For the first time that $P$ touches the line $y=x+1$ at point $(a,a+1)$ ($0 \leq a \leq m - 1$), we have taken $a$ RIGHT steps, and the number of UP steps we have taken so far is $a+1$. From this point, we still need to take $m - a$ RIGHT steps and $n-a-1$ UP steps to reach $(m,n)$. If we transpose this sub-path from $(a,a+1)$ to $(m,n)$ by switching UP step to RIGHT step and RIGHT step to UP step, we have a new lattice path called $Q$. This path $Q$ is a lattice path from $(0,0)$ to $(n-1, m + 1)$ since we need to take $m-a + a + 1 = m+1$ UP steps and $n - 1- a + a= n - 1$ RIGHT steps.

We consider the set $D_{m,n}$ of all lattice paths that go from $(0,0)$ to $(n-1, m + 1)$ by only moving UP or RIGHT for each step. Obviously, we have $|D_{m,n}| = \binom{m+1+n-1}{n-1}=\binom{m+n}{n-1}$.

Following the steps above can map a path in $B_{m,n}$ to a path in $D_{m,n}$.

Let $M \in D_{m,n}$ be a path that goes from $(0,0)$ to $(n-1, m + 1)$ by only moving UP or RIGHT for each step. We can conclude that the number of UP steps in the path $M$
must be at least $2$ steps larger than the number of RIGHT steps ($m+1-(n-1)=m -n + 2 \geq 2$). Therefore, this path must touch the line $y=x+1$ at least once at the point $(a,a+1)$. At this point, we have the number of RIGHT steps that we have taken so far as $a$, and the number of UP steps that we have taken as $a+1$. From this point, we still need to take $n-1-a$ RIGHT steps and $m+1-a-1=m-a$ UP steps. If we transpose this sub-path from $(a,a+1)$ to $(n-1,m+1)$ by switching UP step to RIGHT step and RIGHT step to UP step, we have a new lattice path called $N$. This path $N$ is a lattice path from $(0,0)$ to $(m, n)$ that touches the line $y=x+1$ once and goes above the line $y=x$. 

Following the steps above can map a path in $D_{m,n}$ to a path in $B_{m,n}$.

Therefore, we can create a bijection $f:B_{m,n} \mapsto D_{m,n}$. So $|D_{m,n}|=|B_{m,n}|=\binom{m+n}{n-1}$. Therefore, we can conclude:

$$|\mathcal{C}_{m,n}|=|A_{m,n}| - |B_{m,n}| = \binom{m+n}{n} - \binom{m+n}{n-1}$$

    \item \textbf{Case 2:} $m < n $
    
    Obviously, we have no path in $A_{m,n}$ that will not go over the line $y = x$ since we need to take more UP steps than RIGHT steps. In this case, $|A_{m,n}|=|B_{m,n}|$. Therefore, $|\mathcal{C}_{m,n}|=|A_{m,n}|-|B_{m,n}|=0$. 
  
\end{itemize}

Therefore, with $m < n$, we have $\mathcal{C}_{m,n}=0$ and with $m \geq n$, we have $\mathcal{C}_{m,n}=\binom{m+n}{n} - \binom{m+n}{n-1}$.
 
\vskip .75 in
\item
\textbf{Suppose that $B_n$ is the set of all sequences of $n$ integers $\left\{a_1,a_2,...,a_n\right\}$ where}
\begin{center}$1 \leq a_1 \leq a_2 \leq \cdots \leq a_n\leq n$ and $a_i \leq i$ for all $i$.\end{center}

\textbf{Find (with proof) a bijection between $\mathcal{C}_n$ (see \#1) and $B_n$, thus proving $|B_n| = C_n$.}

We need to find the bijection $f: \mathcal{C}_n \mapsto B_n$ and the inverse function $g: B_n \mapsto \mathcal{C}_n $.

\begin{proof}
    Let $P \in \mathcal{C}_n $ be a lattice path from $(0,0)$ to $(n,n)$ that does not go above the line $y=x$. There are $n$ UP steps and $n$ RIGHT steps in the path.
    
    For each $i, 1 \leq i \leq n$, we have $y_i$ as the $y-$coordinate of the path at the time the $i^{th}$ right step is taken. Since $P$ is a valid path, we must have $y_i \leq i - 1$ since the $y-$coordinate must always be smaller or equal to the $x-$coordinate (For example, when we take the RIGHT step to go from $(2,y)$ to $(3,y)$ or at the time the $3^{rd}$ step is taken, we must have $y \leq 2 = 3 - 1$). Since the lattice path only moves UP or RIGHT, its $y-$coordinate will not decrease (which means that the lattice path never goes down). Obviously, the $y-$coordinate always starts from $0$. Therefore, we have:
    
    $$0 \leq y_1 \leq y_2 \leq y_3 \leq \ldots \leq y_n \leq n - 1$$
    
    With $1 \leq i \leq n$, we define the function $f$ such that $f: \mathcal{C}_n  \mapsto B_n$, $f(P) = A$, $A = \{a_1,a_2,\ldots, a_n\}$, and $a_i = y_i+1$. From this, we have $a_i=y_i+1\leq i - 1 + 1 = 1$. Then with $a_i \leq i$, we have:
    
    $$1 \leq a_1 \leq a_2 \leq a_3 \leq \ldots \leq a_n \leq n$$

    Therefore, we can conclude that $A$ always belongs to the set $B$ since $A$ is the sequence of $n$ integers where $1 \leq a_1 \leq a_2 \leq a_3 \leq \ldots \leq a_n \leq n$. Therefore, we have with every $P \in \mathcal{C}_n$, there exists $A \in B_n$ such that $h(P) = A$.

    Let $A = \{a_1,a_2,a_3,\ldots,a_n\} \in B_n$. For each $a_i \in A$, we declare $y_i = a_1$. We define the function $g:B_n \mapsto \mathcal{C}_n$ such that $g(A) = P$ where $P \in \mathcal{C}_n$. Starting from $(0,0)$, for each $i$, $1 \leq i \leq n$, we do:

\begin{itemize}
    \item We go from the current $x-$coordinate to $i$ .
    \item Let $y_i = a_i-1$. We move UP until the $y-$coordinate reaches $y_i$ (if we are already at $y_i$, take no UP steps)
\end{itemize}

Because $1 \leq a_1 \leq a_2 \leq a_3 \leq \ldots \leq a_n \leq n$, we have $0 \leq y_1 \leq y_2 \leq y_3 \leq \ldots \leq y_n \leq n-1$. This makes sure that this path will always start from $0$ and never go down. Since $i$ is increasing accordingly, we can conclude that the path only has UP and RIGHT steps. There are $n$ RIGHT steps (because we move the $x-$coordinate to the right for $n$ times) and $n$ UP steps (because we need to move the $y-$coordinate from $0$ to $n-1$). Therefore, this path is going from $(0,0)$ to $(n,n)$.

At $x = i-1$, before taking the $i$-th RIGHT step, our $y-$coordinate is $y_i$. Because $A \in B_n$, we know $a_i \leq i$, which means $y_i = a_i - 1 \leq i - 1$. Therefore, the coordinate $(i-1, y_i)$ satisfies $y_i$ smaller or equal to the $x-$coordinate, which means the path never goes above the line $y=x$.

Therefore, we have with every $A \in B_n$, there exists $P \in C_n$ such that $g(A) = P$.

 If we take a path $P \in \mathcal{C}_n$ and let $y_i$ be the $y$-coordinate of the path exactly when the $i$-th RIGHT step is taken, applying the function $f$ maps these coordinates to a sequence $A = \{a_1, a_2, \ldots, a_n\}$ where $a_i = y_i + 1$. Then applying the function $g$ to this sequence $A$ gets the original $y$-coordinates via $y_i = a_i - 1$. By applying these $y_i$ values for each corresponding $i$-th RIGHT step, we can reconstruct the original lattice path $P$. Therefore, $g(f(P)) = P$.
 
 We can apply $g$ to a valid sequence $A$ to construct a Catalan path. If we then extract the $y$-coordinates of the RIGHT steps from this reconstructed path using $f$, we can retrieve the original terms $a_i$, which means $f(g(A)) = A$.

Since $f$ and $g$ are well-defined inverses of each other, $f$ is a bijection between $\mathcal{C}_n$ and $B_n$. Therefore, $|B_n| = |\mathcal{C}_n| = C_n$.

\end{proof}




\vskip .75 in


\item
\begin{enumerate}

\item
\textbf{Let $J_n$ be the set of ways to make $n$ disjoint pairs out of $2n$ points (evenly spaced around a circle) by using $n$ non-intersecting chords and let $j_n=|J_n|$.  Here are the five elements of $J_3$:}\begin{center}
\begin{tikzpicture}[scale=.25]

\def \n {6} 
\def \radius {2.5cm}
\def \margin {8} 
\begin{scope}[color=black,line width=1pt]  

\foreach \s in {1,...,\n}
{
  \node[draw, circle,scale=.5] at ({360/\n * (\s - 1)}:\radius) (v\s) {};  
}

\path
	(v1) edge (v2) 
	(v3) edge (v4)
	(v5) edge (v6);
		
\end{scope}
\end{tikzpicture} 
\quad\quad
\begin{tikzpicture}[scale=.25]
\def \n {6} 
\def \radius {2.5cm}
\def \margin {8} 
\begin{scope}[color=black,line width=1pt]  

\foreach \s in {1,...,\n}
{
  \node[draw, circle,scale=.5] at ({360/\n * (\s - 1)}:\radius) (v\s) {};  
}

\path
	(v1) edge (v2) 
	(v3) edge (v6)
	(v5) edge (v4);
		
\end{scope}
\end{tikzpicture}
\quad\quad
\begin{tikzpicture}[scale=.25]
\def \n {6} 
\def \radius {2.5cm}
\def \margin {8} 
\begin{scope}[color=black,line width=1pt]  

\foreach \s in {1,...,\n}
{
  \node[draw, circle,scale=.5] at ({360/\n * (\s - 1)}:\radius) (v\s) {};  
}

\path
	(v1) edge (v4) 
	(v3) edge (v2)
	(v5) edge (v6);
		
\end{scope}
\end{tikzpicture}
\quad\quad
\begin{tikzpicture}[scale=.25]
\def \n {6} 
\def \radius {2.5cm}
\def \margin {8} 
\begin{scope}[color=black,line width=1pt]  

\foreach \s in {1,...,\n}
{
  \node[draw, circle,scale=.5] at ({360/\n * (\s - 1)}:\radius) (v\s) {};  
}

\path
	(v1) edge (v6) 
	(v3) edge (v2)
	(v5) edge (v4);
		
\end{scope}
\end{tikzpicture}
\quad\quad
\begin{tikzpicture}[scale=.25]
\def \n {6} 
\def \radius {2.5cm}
\def \margin {8} 
\begin{scope}[color=black,line width=1pt]  

\foreach \s in {1,...,\n}
{
  \node[draw, circle,scale=.5] at ({360/\n * (\s - 1)}:\radius) (v\s) {};  
}

\path
	(v1) edge (v6) 
	(v5) edge (v2)
	(v3) edge (v4);
		
\end{scope}
\end{tikzpicture}
\end{center}

\textbf{Prove that $j_n$ satisfies the Catalan recurrence relation.}

The Catalan recurrence relation is $C_n= \sum\limits^{n}_{k=1}C_{k-1}C_{n-k}$.

We have $J_n$ as the set of ways to make $n$ disjoint pairs out of $2n$ points (evenly spaced around a circle) by using $n$ non-intersecting chords.

With $0$ point, we have $1$ way to make $0$ disjoint pairs or $j_0 = 1$. With $2$ points, we have $1$ way to make $1$ disjoin pair or $j_1=1$.

We have $2n$ points evenly spaced around a circle ($n \in \mathbb{N}, 1 \leq n$). We can label the points $1,2, \ldots,2n$ in clockwise order. We start with the point $1$ and connect this point with point $m, m\neq 1$. 
\begin{center}
\begin{tikzpicture}[scale=.25]

\def \n {6} 
\def \radius {2.5cm}
\def \margin {8} 
\begin{scope}[color=black,line width=1pt]  

\foreach \s in {1,...,\n}
{
  \node[draw, circle,scale=.5] at ({360/\n * (\s - 1)}:\radius) (v\s) {};  
}

\node at (3,2.5) {1};

\end{scope}
\end{tikzpicture} 
\quad\quad
\end{center}

Because the chords cannot intersect with each other, the chord between $1$ and $m$ must divide the circle into two disjoint subgroups of points. Let $A$ be the subgroup containing the remaining $2k-2$ points ($1 \leq k \leq n$), the point $1$, and $m$. Therefore, there are $2k$ points in $A$. The subgroup $B$ contains $2n - (2k-2) - 2 = 2(n-k)$ points. 

\begin{center}
\begin{tikzpicture}[scale=.25]

\def \n {6} 
\def \radius {2.5cm}
\def \margin {8} 
\begin{scope}[color=black,line width=1pt]  

\foreach \s in {1,...,\n}
{
  \node[draw, circle,scale=.5] at ({360/\n * (\s - 1)}:\radius) (v\s) {};  
}

\node at (-4,0) {A};
\node at (4,0) {B};
\node at (3,2.5) {1};
\node at (-3,-2.5) {m};
\path
	(v2) edge (v5);	
\end{scope}
\end{tikzpicture} 
\quad\quad
\end{center}

Now, in subgroup $A$, we already have a pair of points connected. We still have $2k-2=2(k-1)$ points waiting to be connected by non-intersecting chords. The number of ways to do this is  the number of way to make $k-1$ disjoint pairs out of $2(k-1)$, or $j_{k-1}$. In subgroup $B$, we have $2n-2k=2(n-k)$ points left waiting to be connected by non-intersecting chords. The number of ways to do this is the number of way to make $n-k$ disjoint pairs out of $2(n-k)$, or $j_{n-k}$. Therefore, the number of ways to connect the remaining $2n-2$ points by splitting them into subgroups $A$ and $B$ is $j_{n-k}j_{k-1}$.

For the first step, there are many options for us to choose the suitable pair of points to connect to create the subgroup $A$ since $1\leq k \leq n$. Therefore, the number of ways to make $n$ disjoint pairs out of $2n$ points (evenly spaced around a circle) by using $n$ non-intersecting chords

$$j_n=\sum\limits_{k=1}^nj_{n-k}j_{k-1}.$$

This recurrence relation looks similar to the Catalan recurrence relation. Therefore, we can conclude that $j_n$ satisfies the Catalan recurrence relation.
\item
\textbf{Let $K_n$ be the set of ways of pairing $2n$ points laying on a horizontal line using non-intersecting arcs that are drawn above the points and let $k_n=|K_n|$.  Here are the five elements of $K_3$:}\begin{center}

\begin{tikzpicture}[scale=.25]  
  \begin{scope}[color=black,line width=1pt]  %
    	
	\path (2,0) node[draw,shape=circle,scale=.5] (v1) {};  
	 \path (4,0) node[draw,shape=circle,scale=.5] (v2) {};
	\path (6,0) node[draw,shape=circle,scale=.5] (v3) {};
	\path (8,0) node[draw,shape=circle,scale=.5] (v4) {};
	\path (10,0) node[draw,shape=circle,scale=.5] (v5) {};
	\path (12,0) node[draw,shape=circle,scale=.5] (v6) {};
	
	\path
	(v1) edge[out=60,in=120] (v2)
	(v3) edge[out=60,in=120] (v4)
	(v5) edge[out=60,in=120] (v6);
		
\end{scope}
\end{tikzpicture}
\quad\quad
\begin{tikzpicture}[scale=.25]  
  \begin{scope}[color=black,line width=1pt]  %
    	
	\path (2,0) node[draw,shape=circle,scale=.5] (v1) {};  
	 \path (4,0) node[draw,shape=circle,scale=.5] (v2) {};
	\path (6,0) node[draw,shape=circle,scale=.5] (v3) {};
	\path (8,0) node[draw,shape=circle,scale=.5] (v4) {};
	\path (10,0) node[draw,shape=circle,scale=.5] (v5) {};
	\path (12,0) node[draw,shape=circle,scale=.5] (v6) {};
	
	\path
	(v1) edge[out=60,in=120] (v2)
	(v3) edge[out=60,in=120] (v6)
	(v4) edge[out=60,in=120] (v5);
		
\end{scope}
\end{tikzpicture}
\quad\quad
\begin{tikzpicture}[scale=.25]  
  \begin{scope}[color=black,line width=1pt]  %
    	
	\path (2,0) node[draw,shape=circle,scale=.5] (v1) {};  
	 \path (4,0) node[draw,shape=circle,scale=.5] (v2) {};
	\path (6,0) node[draw,shape=circle,scale=.5] (v3) {};
	\path (8,0) node[draw,shape=circle,scale=.5] (v4) {};
	\path (10,0) node[draw,shape=circle,scale=.5] (v5) {};
	\path (12,0) node[draw,shape=circle,scale=.5] (v6) {};
	
	\path
	(v1) edge[out=60,in=120] (v6)
	(v2) edge[out=60,in=120] (v3)
	(v4) edge[out=60,in=120] (v5);
		
\end{scope}
\end{tikzpicture}
\end{center}
\begin{center}
\begin{tikzpicture}[scale=.25]  
  \begin{scope}[color=black,line width=1pt]  %
    	
	\path (2,0) node[draw,shape=circle,scale=.5] (v1) {};  
	 \path (4,0) node[draw,shape=circle,scale=.5] (v2) {};
	\path (6,0) node[draw,shape=circle,scale=.5] (v3) {};
	\path (8,0) node[draw,shape=circle,scale=.5] (v4) {};
	\path (10,0) node[draw,shape=circle,scale=.5] (v5) {};
	\path (12,0) node[draw,shape=circle,scale=.5] (v6) {};
	
	\path
	(v1) edge[out=60,in=120] (v4)
	(v2) edge[out=60,in=120] (v3)
	(v5) edge[out=60,in=120] (v6);
		
\end{scope}
\end{tikzpicture}
\quad\quad
\begin{tikzpicture}[scale=.25]  
  \begin{scope}[color=black,line width=1pt]  %
    	
	\path (2,0) node[draw,shape=circle,scale=.5] (v1) {};  
	 \path (4,0) node[draw,shape=circle,scale=.5] (v2) {};
	\path (6,0) node[draw,shape=circle,scale=.5] (v3) {};
	\path (8,0) node[draw,shape=circle,scale=.5] (v4) {};
	\path (10,0) node[draw,shape=circle,scale=.5] (v5) {};
	\path (12,0) node[draw,shape=circle,scale=.5] (v6) {};
	
	\path
	(v1) edge[out=60,in=120] (v6)
	(v2) edge[out=60,in=120] (v5)
	(v3) edge[out=60,in=120] (v4);
		
\end{scope}
\end{tikzpicture}

\end{center}
\textbf{Find (with proof) a bijection between $K_n$ and $J_n$, thus proving $|K_n| = C_n$.}

The Catalan recurrence relation is $C_n= \sum\limits^{n}_{k=1}C_{k-1}C_{n-k}$.

We need to define two functions: 
\begin{itemize}
    \item $f:J_n \mapsto K_n$ such that for every $J \in J_n$, there exists $K \in K_n$ such that $f(J) =K$
    \item $g:K_n \mapsto J_n$ such that for every $K \in K_n$, there exists $J \in J_n$ such that $g(K) =J$
\end{itemize}

We have $K_n$ be the set of ways of pairing $2n$ points laying on a horizontal line using non-intersecting arcs that are drawn above the points. We can label the points $1,2, \ldots,2n$ from left to right. 

We have $J_n$ as the set of ways to make $n$ disjoint pairs out of $2n$ points (evenly spaced around a circle) by using $n$ non-intersecting chords. We can label the points $1,2, \ldots,2n$ in clockwise order. 

We already know that in the circle, two chords connecting pairs $(x_1, y_1)$ and $(x_2, y_2)$  (Note: $x_1 < y_1$ and  $x_2 < y_2$) intersects if and only if $x_1 < x_2 < y_1 < y_2$.

On the horizontal line, two overhead arcs connecting $(x_1, y_1)$ and $(x_2, y_2)$ (Note: $x_1 < y_1$ and  $x_2 < y_2$) cross each other if and only if $x_1 < x_2 < y_1 < y_2$.

We first define function $f:K_n \mapsto J_n$. We have for a $K \in K_n$, if there is an arc connecting point $i$ and point $j$ (where $i < j$) on the horizontal line, we define $f(K)$ by drawing a chord between two points $i$ and $j$ in the evenly spaced circle. 

Because $K \in K_n$ only has non-intersecting arcs, there exist no two arcs that have endpoints $(x_1, y_1)$ and $(x_2, y_2)$ such that $x_1 < x_2 < y_1 < y_2$. Therefore, there will not exist any two corresponding chords  $(x_1, y_1)$ and $(x_2, y_2)$ in $f(K)$ that have $x_1 < x_2 < y_1 < y_2$, or there are no pairs of chords in $f(K)$ that are intersecting. Therefore, $f(K)$ belongs to $J_n$.

We then define function $g:J_n \mapsto K_n$. We have for a $J \in J_n$, if there is a chord connecting point $i$ and point $j$ (where $i < j$) on the circle, we define $g(J)$ by drawing an overhead arc between two points $i$ and $j$ on the horizontal line. 

Because $J \in J_n$ only has non-intersecting chords, there exist no two chords that have endpoints $(x_1, y_1)$ and $(x_2, y_2)$ such that $x_1 < x_2 < y_1 < y_2$. Therefore, there will not exist any two corresponding arcs that have the corresponding endpoints $(x_1, y_1)$ and $(x_2, y_2)$ in $g(J)$ that have $x_1 < x_2 < y_1 < y_2$, or there are no pairs of arcs in $g(J)$ that are intersecting. Therefore, $g(J)$ belongs to $K_n$.

Because $f$ maps arcs to chords using the exact same pairs of points, applying $f$ and then $g$ returns the original set of arcs ($g(f(K)) = K$). Because $g$ maps chords to arcs using the exact same pairs of points, applying $f$ and then $g$ returns the original set of chords ($f(g(J)) = J$).

Since $f$ and $g$ are well-defined inverses of each other, $f$ is a bijection. Therefore, 
we can conclude that $|K_n| = |J_n|=C_n$.

\end{enumerate}
\vskip 2cm


{\bf Go to next page (posted later)}
\newpage

\item
\begin{enumerate}
\item
\textbf{In class, we saw that the formula for the number of partitions of a non-negative integer $n$ is quite complicated.  Find and prove a formula for the number of partitions of $n$ whose parts are only equal to $2$ or $1$.}

For every $n \in \mathbb{N}$, $a_n$ is the number of partitions of $n$ whose parts are only equal to $1$ or $2$.

We first consider the generating function for the part equal to $1$.
This sum created by the parts that are equal to $1$ can start from 0 and increase by 1. We represent these choices as the exponents of our variable $x$. The generating function for the first integer is the sum of all these possibilities:

\begin{align*}
f_{1} &= x^0+x^1 + x^{1+1} + x^{1+1+1} + \ldots  \\
    &= \sum\limits^{\infty}_{n=0}x^n\\
    &= \frac{1}{1-x} \\
\end{align*}

We then consider the generating function for the part equal to $2$.
This sum, created by the parts that are equal to $2$ can start from 0 and increase by 2. We represent these choices as the exponents of our variable $x$. The generating function for the first integer is the sum of all these possibilities:

\begin{align*}
f_{2} &= x^0+x^2 + x^{2+2} + x^{2+2+2} + \ldots  \\
    &= \sum\limits^{\infty}_{n=0}x^{2n}\\
    &= \frac{1}{1-x^2} \\
\end{align*}

From these, we have the generating function for $\{a_n\}$ by multiplying the generating functions together, because multiplying terms from each series adds their exponents, which counts every valid partition of $n$: 

\begin{align*}
    f_{1} \times f_{2}&= \frac{1}{(1-x^2)(1-x)}\\
    &= \frac{1}{(1-x)^2}\cdot\frac{1}{1+x}\\
\end{align*}

I used Mathematica to split the generating function into factors (Code: $Apart[1/((1-x)^2*(1+x))]$), then we have:
\begin{align*}
    \frac{1}{(1-x)^2}\cdot\frac{1}{1+x} &= \frac{1}{2(1-x)^2} + \frac{1}{4(1-x)} + \frac{1}{4(1+x)}\\
    &=\frac{1}{2}\sum\limits^{\infty}_{n=0}\binom{n+1}{1}x^{n}+\frac{1}{4}\sum\limits^{\infty}_{n=0}x^{n}+\frac{1}{4}\sum\limits^{\infty}_{n=0}(-1)^nx^{n}\\
    &= \sum\limits^{\infty}_{n=0}\left(\frac{1}{2}(n+1)+\frac{1}{4}+\frac{(-1)^n}{4}\right)x^n
\end{align*}

The number of partitions of $n$ with the parts equal to $1$ or $2$ is the coefficient of $x^n$ in the generating function above, or $a_n = \frac{1}{2}(n+1)+\frac{1}{4}+\frac{(-1)^n}{4}$.

We consider the parity of $n$:
\begin{itemize}
    \item \textbf{Case 1:} $n \equiv 0 \pmod 2$, so we can write $n = 2k, k \in \mathbb{N}$
    
    We have $a_n = \frac{1}{2}(n+1)+\frac{1}{4}+\frac{1}{4} = \frac{n}{2} + 1 = k+1$.

    \item \textbf{Case 2:} $n \equiv 1 \pmod 2$,  so we can write $n = 2k+1, k \in \mathbb{N}$
    
    We have $a_n = \frac{1}{2}(n+1)+\frac{1}{4}-\frac{1}{4} = \frac{n+1}{2}=\frac{2k+2}{2} = k+1.$
\end{itemize}

\item
\textbf{In class, we saw that the number of compositions of a non-negative integer $n$ is $2^{n-1}$ (much nicer-looking than the partitions formula).  Find and prove a formula for the number of compositions of $n$ whose parts are only $1$'s and $2$'s.}


For every $n \in \mathbb{N}$, $b_n$ is the number of compositions of $n$ whose parts are only $1$'s and $2$'s.

With $n = 0$, we have $b_0 = 1$ since we can have an empty sequence that satisfies the sum.

With $n = 1$, we have $b_1 = 1$, and the only composition that satisfies the sum is $\{1\}$.

If we start the composition with a $1$, we have to find a composition of $n-1$ to add to the sequence so we can get a composition of $n$. There are $b_{n-1}$ compositions of $n-1$ we can use.

If we start the composition with a $2$, we have to find a composition of $n-2$ to add to the sequence so we can get a composition of $n$. There are $b_{n-2}$ compositions of $n-2$ we can use.

Therefore, the total number of compositions of $n$ whose parts are only $1$'s and $2$'s is 

$$b_0=1, b_1=1, b_n = b_{n-1}+b_{n-2}.$$

We have the characteristic equation: $x^2 - x - 1 = 0$. We can easily find two roots of this equation: $\lambda_1 = \frac{1+\sqrt{5}}{2}, \lambda_2 = \frac{1-\sqrt{5}}{2}$. Since we have $\lambda_1 \neq \lambda_2$, we can conclude that $a_n=c_1\cdot \lambda_1^n+c_2\cdot \lambda_2^n$.

We need to find $c_1$ and $c_2$. Let $n = 0,1$, we have:

$$\begin{cases}
     c_1 \cdot \lambda_1^0 + c_2 \cdot \lambda_2^0 &= b_0 \\
    c_1 \cdot \lambda_1^1 + c_2 \cdot \lambda_2^1 &= b_1
\end{cases}$$

Then we substitute the value of $\lambda_1$ and $\lambda_2$:

$$\begin{cases}
    c_1 + c_2 &= 1\\
    \frac{1+\sqrt{5}}{2}c_1  + \frac{1-\sqrt{5}}{2}c_2 &= 1
\end{cases}$$

Solving these two equations, we have $c_1 = \frac{1}{\sqrt{5}} \cdot \left(\frac{1+\sqrt{5}}{2}\right)$ and $c_2 = \frac{-1}{\sqrt{5}} \cdot \left(\frac{1-\sqrt{5}}{2}\right)$. Therefore, the exact formula for $a_n$ is:

$$b_n=\frac{1}{\sqrt{5}} \cdot \left(\frac{1+\sqrt{5}}{2}\right)^{n+1}-\frac{1}{\sqrt{5}} \cdot \left(\frac{1-\sqrt{5}}{2}\right)^{n+1}$$
\end{enumerate}
\vskip 2cm



\item
\textbf{Define $p'(n,k)$ to be the number of partitions of $n$ whose largest part is equal to $k$.}
\begin{enumerate}
\item \textbf{Make a triangle (analogous to Pascal's triangle) where the $k$-th entry in the $n$-th row equals $p'(n,k)$. You should make one with at least 8 rows total.  Remember to start with ``row 0" and each row should begin with ``entry 0."}
\begin{table}[htbp]
\centering
\renewcommand{\arraystretch}{1.2}
    \begin{tabular}{c | c c c c c c c c c}
    $n \setminus k$ & \textbf{0} & \textbf{1} & \textbf{2} & \textbf{3} & \textbf{4} & \textbf{5} & \textbf{6} & \textbf{7} & \textbf{8} \\
    \hline
    \textbf{0} & 1 & & & & & & & & \\
    \textbf{1} & 0 & 1 & & & & & & & \\
    \textbf{2} & 0 & 1 & 1 & & & & & & \\
    \textbf{3} & 0 & 1 & 1 & 1 & & & & & \\
    \textbf{4} & 0 & 1 & 2 & 1 & 1 & & & & \\
    \textbf{5} & 0 & 1 & 2 & 2 & 1 & 1 & & & \\
    \textbf{6} & 0 & 1 & 3 & 3 & 2 & 1 & 1 & & \\
    \textbf{7} & 0 & 1 & 3 & 4 & 3 & 2 & 1 & 1 & \\
    \textbf{8} & 0 & 1 & 4 & 5 & 5 & 3 & 2 & 1 & 1 \\
    \end{tabular}
\end{table}
\item
\textbf{Prove that $p'(n,k)=p'(n-1,k-1)+p'(n-k,k)$ for $0<k<n$.  You may assume $p'(n,0)=0$ and $p'(0,0)=1$.}

We have $p'(n,k)$, which is the number of partitions of $n$ whose largest part is equal to $k$, $p'(n-1,k-1)$, which is the number of partitions of $n-1$ whose largest part is equal to $k-1$ and $p'(n-k,k)$, which is the number of partitions of $n-k$ whose largest part is equal to $k$.

Let $P'_{n,k}$ be the set of partitions of $n$ that has the largest part equal to $k$ ($|P'_{n,k}|=p'(n,k)$). Let $P'_{n-1,k-1}$ be the set of partitions of $n-1$ that has the largest part equal to $k-1$ ($|P'_{n-1,k-1}|=p'(n-1,k-1)$). Let $P'_{n-k,k}$ be the set of partitions of $n-k$ that has the largest part equal to $k$ ($|P'_{n-k,k}|=p'(n-k,k)$).

We can divide the set $P'_{n,k}$ into two mutually exclusive and disjoint subsets:
\begin{itemize}
    \item \textbf{Subset $P^1_{n,k}$:} The partitions where the largest part $k$ appears exactly once.
    
    \item \textbf{Subset $P^2_{n,k}$:} The partitions where the largest part $k$ appears at least twice.
\end{itemize}

With a partition $A \in P'_{n-1,k-1}$, we have a sequence of non-increasing positive integers that sum up to $n-1$ and there exists at least one part that is equal to $k-1$. We can add $1$ to one (and only one) of the parts $k-1$ to this sequence. Now, we have a new sequence $A'$ with the part that has the largest value equal to $k$ (since $k-1 < k$) and the sum of this sequence is $n-1 + 1 = n$. This new sequence $A'$ belongs to $P^1_{n,k}$.

With a partition $B \in P^1_{n,k}$, we have a sequence of non-increasing positive integers that sum up to $n$ and there exists only one part that is equal to $k$. We can deduct $1$ from this part of $k$ in this sequence. Now, we have a new sequence $B'$ with the part that has the largest value equal to $k-1$ and the sum of this sequence is $n-1$. This new sequence $B'$ belongs to $P^1_{n-1,k-1}$.

Therefore, we can find a bijection $f:P^1_{n,k} \mapsto P'_{n-1,k-1}$, or $|P^1_{n,k}|=|P'_{n-1,k-1}|$.

With a partition $C \in P'_{n-k,k}$, we have a sequence of non-increasing positive integers that sum up to $n-k$ and there exists at least one part that is equal to $k$ in the sequence. We can simply add a new part $k$ to this sequence. Therefore, we have a new sequence $C'$ with the part that has the largest value equal to $k$ and the sum of this sequence is $n-k + k = n$. This new sequence $C'$ belongs to $P^2_{n,k}$ since it has at least two parts that are equal to $k$.

With a partition $D \in P^2_{n,k}$, we have a sequence of non-increasing positive integers that sum up to $n$, and there exists at least two parts that are equal to $k$ in the sequence. We can simply remove one part $k$ from this sequence. Therefore, we have a new sequence $D'$ with at least one part that has the largest value and it is equal to $k$. The sum of this sequence is $n-k$. This new sequence $D'$ belongs to $P'_{n-k,k}$ accordingly.

Therefore, we can find a bijection $g:P^2_{n,k} \mapsto P'_{n-k,k}$, or $|P^2_{n,k}|=|P'_{n-k,k}|$.

With every partition $P \in P'_{n-1,k-1}$, we can construct a unique partition that belongs to $P'_{n,k}$ and has only one largest part that is equal to $k$. With every partition $P \in P'_{n-k,k}$, we can construct a unique partition that belongs to $P'_{n,k}$ and has at least two parts that are equal to $k$. Because every partition in $P'_{n,k}$ must belong to exactly one of these two subsets, the total number of partitions $p'(n,k)$ is 

$$|P'_{n,k}| = |P'_{n-1,k-1}|+ |P'_{n-k,k}|$$

Therefore, we can conclude that

$$p'(n,k)=p'(n-1,k-1)+p'(n-k,k).$$


\end{enumerate}
\vskip 2cm

\item
\begin{enumerate}
\item
\textbf{Suppose $A$ and $B$ are finite sets and $f:A\to B$ is a one-to-one function (i.e. an injection) that is also {\it not} onto (i.e. not surjective).  Explain why this proves $|A| < |B|$.}

In the function $f$, we have the domain $A$ and the codomain $B$. 

Since $f:A \to B$ is a one-to-one function, when we apply $f$ onto each element $a \in A$, we get a unique element $b \in B$, or if we have $a, a' \in A$, $a \neq a'$, we ensure that $f(a)=b \neq b' = f(a')$ where $b, b' \in B$. If we have $|A|$ elements in $A$, we can initially conclude that there are at least $|A|$ unique elements in $B$ to satisfy this, or $|B| \geq |A|$.

However, $f:a \to B$ is not an onto function, which means there always exists at least one element $b \in B$ that is not mapped from any element $a$ of the domain $A$. Therefore, $B$ may have at least more than $1$ element compared to $A$, or $|B| \geq |A| + 1 > |A|$.

In conclusion, we have:

$$|A| < |B|$$

\item \textbf{Prove that $p(n-1)<p(n)$ for all $n>1$.}

\begin{proof}
    We have with $n>1$, $p(n-1)$ is the number of partitions of $n-1$, or the number of set of non-increasing list of positive integers that sum to $n-1$. Let $A_{n-1}$ be the set of set of non-increasing list of positive integers that sum to $n-1$, and $A_{n}$ be the set of set of non-increasing list of positive integers that sum to $n$.

    For each list $S \in A_{n-1}$ of non-increasing list of positive integers that sum to $n-1$, if we add $1$ to the $S$, we will have a new list of non-increasing list of positive integers that sum to $n$, or a partition of $n$. If we have two unique lists $S_i, S_j \in A$ such that $i \neq j, S_i \neq S_j$, we ensure that the new lists we create by adding $1$ to them will be different from each other, or $S_i \cup \{1\} \neq S_j \cup \{1\}$. Therefore, we can initially conclude that there are at least $p(n-1)$  unique elements in $p(n)$.
    
    However, the partitions we construct above for $n$ always contain $1$. But there exists the partition $\{n\}$, $n > 1$ that does not contain $1$ but still in $A_n$, which means there exists at least 1 partition in $|A_n|$ that is not constructed from $A_{n-1}$. Therefore, we have $|A_n| \geq |A_{n-1}|+1>|A_{n-1}|$.

    Therefore, with $n > 1$, we have $p(n-1) < p(n)$.
    
\end{proof}


\item \textbf{Prove that $\ds\sum_{k=1}^{n}p(k)<p(2n)$ for all $n>0$.}

\begin{proof}
    Let $P_n$ be the set of partitions of $n$, so we have $|P_n|=p(n)$. 

    We need to define the function $f_k: (P_1 \cup P_2\cup P_3 \cup \ldots P_k) \mapsto P_{2n}$ with $1\leq k \leq n$. The domain is $P_1 \cup P_2\cup P_3 \cup \ldots P_k$, so the size of domain is $\ds\sum_{k=1}^{n}p(k)$. The co-domain is $P_{2n}$, so the size of co-domain is $p(2n)$.
    
    Let $P \in P_k$, $1 \leq k \leq n$. can add a part equal to $2n-k$ to every single partition $P \in P_k$ to construct a new partition $P'$ of $2n -k + k = 2n$, which belongs to $P_{2n}$. In the newly constructed partitions $P' \in p_{2n}$ above, because $k \leq n$, the added part is $2n - k$ is always larger or equal to $n$ ($2n-k \geq n$). Meanwhile, the partition $P$ is partition of $k$, which means the largest possible part in $P$ is $k$. Since $k \leq n$, every part $a_i$ in $P$ must be smaller or equal to $n$ ($a_i \leq n$). Therefore, the added part ($2n - k$) is ensured to be greater than or equal to every part in $P$, which means it will always be the part of the new partition $P'$ that has largest value. 

    For different values of $k$, we will always have unique value of $2n-k$ to add to the existing partitions of different $k$. Every partition in $P_k$ is unique. Therefore, the constructed partition is also unique based on the element from the domain. 

    With the constructed partitions $P'$, we can remove $2n-k$ (the largest part in the partition) to get the partition $P$ of $2n - (2n-k) = k$, which belongs to $P_k$.

    Because we can transform the partition of the domain to a unique partition of in the co-domain and reconstruct some of the partitions of the co-domain into a partition in the domain, $f$ is one-to-one. This proves that $\sum_{k=1}^{n}p(k) \leq p(2n)$.

    Our function $f$ always adds a part of $2n - k$. Because $k \geq 1$, the largest part we can add to is $2n - 1$. Therefore, no partition constructed by the function $f$ can  contain the number $2n$. However, the partition $\{2n\}$ does belong to $P_{2n}$. Since it was not generated by our function, there exists at least one element in the co-domain $P_{2n}$ that is not mapped from any element of the domain. Therefore, $f$ is not onto. 
    
    Therefore, the size of our domain is strictly less than the size of our co-domain. We have:

    $$|P_1 \cup P_2\cup P_3 \cup \ldots P_k| = \ds\sum_{k=1}^{n}p(k) < p(2n) = |P_{2n}|$$
    
\end{proof}
\end{enumerate}

\end{enumerate}

\end{document}
